% impossibility.tex — Standalone LaTeX file; also suitable for \input{} into main paper.
% Proves that no single sequential gradient-boosted model can simultaneously
% satisfy Stability and Equity under high feature collinearity.

\documentclass[11pt]{article}

% ── Packages ──────────────────────────────────────────────────────────────────
\usepackage{amsmath,amssymb,amsthm}
\usepackage{mathtools}
\usepackage[margin=1in]{geometry}
\usepackage{hyperref}
\usepackage{cleveref}

% ── Theorem environments ─────────────────────────────────────────────────────
\newtheorem{theorem}{Theorem}
\newtheorem{lemma}[theorem]{Lemma}
\newtheorem{corollary}[theorem]{Corollary}
\newtheorem{definition}[theorem]{Definition}
\theoremstyle{remark}
\newtheorem{remark}[theorem]{Remark}

% ── Notation shortcuts ───────────────────────────────────────────────────────
\newcommand{\feat}{\mathcal{P}}             % feature set {1,...,P}
\newcommand{\group}{\mathcal{G}}            % feature partition
\newcommand{\Gl}{G_\ell}                    % a single group
\newcommand{\phij}{\phi_j}                  % SHAP attribution for feature j
\newcommand{\phik}{\phi_k}                  % SHAP attribution for feature k
\newcommand{\splits}[1]{n_{#1}}             % number of splits on feature #1
\newcommand{\totalT}{T}                     % total number of boosting rounds
\newcommand{\rhovar}{\rho}                  % pairwise correlation
\newcommand{\Exp}{\mathbb{E}}              % expectation
\newcommand{\Prob}{\mathbb{P}}             % probability
\newcommand{\Var}{\mathrm{Var}}            % variance
\newcommand{\abs}[1]{\lvert #1 \rvert}
\newcommand{\norm}[1]{\lVert #1 \rVert}
\newcommand{\dash}{\textsc{DASH}}

\title{An Impossibility Theorem for Stable and Equitable\\
  Feature Attribution in Sequential Gradient Boosting}
\author{}
\date{}

\begin{document}
\maketitle

% ══════════════════════════════════════════════════════════════════════════════
\section{Formal Setup}
\label{sec:setup}
% ══════════════════════════════════════════════════════════════════════════════

Consider a supervised learning problem with $P$ features
$\feat = \{1, \dots, P\}$ and response $y$.
The features admit a \emph{correlation partition}
$\group = \{\Gl\}_{\ell=1}^{L}$ such that for every group $\Gl$ and every
pair $j, k \in \Gl$, the pairwise Pearson correlation satisfies
$\mathrm{Corr}(x_j, x_k) = \rhovar$ for some $\rhovar \in (0, 1)$.

\begin{definition}[Data-generating process]
\label{def:dgp}
The true DGP is
\begin{equation}
  y = \sum_{\ell=1}^{L} \sum_{j \in \Gl} \beta_j\, x_j + \varepsilon,
  \qquad \varepsilon \sim \mathcal{N}(0, \sigma^2),
\end{equation}
where all features within each group share a common coefficient:
$\beta_j = \beta_k$ for all $j, k \in \Gl$.
\end{definition}

\begin{definition}[Sequential gradient boosting]
\label{def:sgb}
A gradient-boosted ensemble $f$ is constructed sequentially:
\begin{equation}
  f = \sum_{t=1}^{\totalT} \eta\, h_t,
\end{equation}
where $\eta > 0$ is the learning rate and each base learner $h_t$ is a
regression tree fit to the pseudo-residuals
$r_t = y - \sum_{s < t} h_s(x)$.
At each internal node, the tree greedily selects the feature--threshold pair
$(j^*, c^*)$ maximizing the squared-error reduction.
We write $\omega \in \Omega$ for the random seed controlling data
sub-sampling, feature sub-sampling, and tie-breaking during training.
\end{definition}

\begin{definition}[Attributions and metrics]
\label{def:attr}
Let $\phij(f)$ denote the interventional TreeSHAP global feature importance
for feature $j$ in model $f$, defined as the mean absolute SHAP value over
an explanation set:
$\phij(f) = \frac{1}{n} \sum_{i=1}^{n} \abs{\phi_j(f; x^{(i)})}$.
\end{definition}

\begin{definition}[Stability and Equity]
\label{def:desiderata}
Let $f_\omega$ denote a model trained with random seed $\omega$.
\begin{itemize}
  \item \textbf{Stability.} A method is $\delta$-stable if for any two
    independent draws $\omega, \omega'$, the Spearman rank correlation
    between the importance vectors satisfies
    \begin{equation}
      \Exp\bigl[\mathrm{Spearman}\bigl(
        \boldsymbol{\phi}(f_\omega),\;
        \boldsymbol{\phi}(f_{\omega'})
      \bigr)\bigr] \geq 1 - \delta.
    \end{equation}

  \item \textbf{Equity.} A method is $\gamma$-equitable if for every group
    $\Gl$ with $\beta_j = \beta_k$ for all $j, k \in \Gl$,
    \begin{equation}
      \Exp\!\left[
        \frac{\max_{j \in \Gl} \phij(f_\omega)}
             {\min_{k \in \Gl} \phik(f_\omega)}
      \right] \leq 1 + \gamma.
    \end{equation}
\end{itemize}
A method satisfies both desiderata simultaneously if it is
$\delta$-stable and $\gamma$-equitable for small $\delta, \gamma > 0$.
\end{definition}


% ══════════════════════════════════════════════════════════════════════════════
\section{First-Mover Concentration}
\label{sec:first-mover}
% ══════════════════════════════════════════════════════════════════════════════

\begin{lemma}[First-mover concentration of splits]
\label{lem:split-concentration}
Consider a group $\Gl = \{j_1, \dots, j_m\}$ with pairwise correlation
$\rhovar$ and equal true coefficients.  In a sequential gradient-boosted
ensemble $f_\omega$ with $\totalT$ trees, suppose the first tree $h_1$
selects feature $j_1 \in \Gl$ at its root split (determined by seed
$\omega$).  Then for any other feature $j_q \in \Gl$, $q \neq 1$, the
expected cumulative split counts satisfy
\begin{equation}
  \Exp_\omega\bigl[\splits{j_1}\bigr]
    - \Exp_\omega\bigl[\splits{j_q}\bigr]
    = \Omega(\rhovar \cdot \totalT).
  \label{eq:split-gap}
\end{equation}
\end{lemma}

\begin{proof}
Because $\mathrm{Corr}(x_{j_1}, x_{j_q}) = \rhovar$, once $h_1$ splits on
$j_1$ and absorbs a fraction of the signal shared with $j_q$, the
pseudo-residual $r_2 = y - \eta\, h_1(x)$ retains the component of $y$
orthogonal to $x_{j_1}$.  The residual variance explained by $j_q$ in
subsequent trees is reduced by a factor proportional to $\rhovar^2$.

Formally, write $x_{j_q} = \rhovar\, x_{j_1} + \sqrt{1-\rhovar^2}\, z_q$
where $z_q \perp x_{j_1}$.  The portion of the signal along $x_{j_1}$
has already been captured by $h_1$; the marginal squared-error reduction
from splitting on $j_q$ in tree $h_2$ is proportional to
$(1 - \rhovar^2)\, \beta_{j_q}^2\, \Var(z_q)$,
while the marginal gain from a further split on $j_1$ captures
higher-order residual structure at full variance $\Var(x_{j_1})$.

Thus $j_1$ has a \emph{per-round advantage} of magnitude
$\Theta(\rhovar^2)$ in expected gain.
Since this advantage persists across all $\totalT$ rounds (the residuals
continue to align preferentially with $j_1$), the cumulative split-count
gap grows as
\begin{equation}
  \Exp[\splits{j_1}] - \Exp[\splits{j_q}]
    \;\geq\; c \cdot \rhovar^2 \cdot \totalT
\end{equation}
for a constant $c > 0$ that depends on $\eta$, the tree depth, and
$\beta_{j_1}$.  Since $\rhovar^2 = \Theta(\rhovar)$ for
$\rhovar \to 1$, this establishes the $\Omega(\rhovar \cdot \totalT)$
bound.

Crucially, the identity of the first-mover $j_1$ is determined by
$\omega$ (through sub-sampling and tie-breaking) and is
\emph{approximately uniform} over $\Gl$ by symmetry of the DGP.
\end{proof}


% ══════════════════════════════════════════════════════════════════════════════
\section{Attribution Inheritance}
\label{sec:attribution}
% ══════════════════════════════════════════════════════════════════════════════

\begin{lemma}[TreeSHAP inherits split concentration]
\label{lem:shap-inherits}
For a tree ensemble $f = \sum_{t} h_t$, the interventional TreeSHAP
global importance $\phij(f)$ satisfies
\begin{equation}
  \phij(f) = \Theta\!\left(
    \frac{\splits{j}}{\sum_{k=1}^{P} \splits{k}}
  \right).
  \label{eq:shap-splits}
\end{equation}
Consequently, if $\splits{j_1} - \splits{j_q} = \Omega(\rhovar \cdot
\totalT)$, then
\begin{equation}
  \phi_{j_1}(f) - \phi_{j_q}(f) = \Omega(\rhovar).
  \label{eq:shap-gap}
\end{equation}
\end{lemma}

\begin{proof}
Interventional TreeSHAP computes, for each tree $h_t$, the Shapley value
of feature $j$ with respect to the prediction function defined by that
tree's structure.  A feature that never appears as a split variable in
$h_t$ receives zero attribution from that tree.  For a feature that does
appear, its attribution in $h_t$ is determined by the magnitude of the
prediction difference at the nodes where it splits, weighted by the
Shapley combinatorial coefficients.

Summing over all trees, the global importance $\phij(f)$ aggregates
contributions from exactly those trees containing splits on $j$.
Under standard regularity conditions (bounded tree depth $d$, learning
rate $\eta$, and feature values), each split on $j$ contributes
$\Theta(\eta)$ to $\phij(f)$ in expectation.
Therefore
\begin{equation}
  \phij(f)
    = \Theta(\eta \cdot \splits{j})
    = \Theta\!\left(\frac{\splits{j}}{\sum_k \splits{k}}\right),
\end{equation}
since $\sum_k \splits{k} = \Theta(\totalT)$ for trees of bounded depth.

Applying \Cref{lem:split-concentration}: the split gap
$\splits{j_1} - \splits{j_q} = \Omega(\rhovar \cdot \totalT)$ implies
\begin{equation}
  \phi_{j_1}(f) - \phi_{j_q}(f)
    = \Theta\!\left(\eta \cdot \rhovar \cdot \totalT
      \cdot \frac{1}{\totalT}\right)
    = \Omega(\rhovar),
\end{equation}
which is bounded away from zero for any fixed $\rhovar > 0$.
\end{proof}


% ══════════════════════════════════════════════════════════════════════════════
\section{Impossibility Theorem}
\label{sec:impossibility}
% ══════════════════════════════════════════════════════════════════════════════

\begin{theorem}[Impossibility of simultaneous Stability and Equity]
\label{thm:impossibility}
Let $\Gl = \{j_1, \dots, j_m\}$ be a group of $m \geq 2$ features with
pairwise correlation $\rhovar$ and equal true coefficients.
For any sequential gradient-boosted ensemble $f_\omega$ with $\totalT$
trees and any $\rhovar \in (0,1)$:
\begin{enumerate}
  \item[(i)] $\gamma$-Equity requires $\gamma \to \infty$ as
    $\rhovar \to 1$, i.e., the attribution ratio within $\Gl$ diverges; or
  \item[(ii)] if Equity is enforced (e.g., by post-hoc normalization), then
    $\delta$-Stability requires $\delta \to 1$, i.e., rankings become
    maximally unstable.
\end{enumerate}
No choice of $\totalT$, learning rate $\eta$, or tree structure
eliminates this trade-off.
\end{theorem}

\begin{proof}
We prove both parts separately.

\medskip
\noindent\textbf{Part (i): Sequential boosting violates Equity.}

Fix a seed $\omega$.  By symmetry of the DGP
(\Cref{def:dgp}), every feature in $\Gl$ has the same marginal
distribution and the same true coefficient.  However, the greedy split
selection in $h_1$ must break this symmetry: exactly one feature
$j^*(\omega) \in \Gl$ is selected at the root of $h_1$, determined by
the random seed $\omega$ through sub-sampling and numerical tie-breaking.

By \Cref{lem:split-concentration}, this first-mover $j^*(\omega)$
accumulates a split advantage of $\Omega(\rhovar \cdot \totalT)$ over
every other member of $\Gl$.  By \Cref{lem:shap-inherits}, this
translates to an attribution gap:
\begin{equation}
  \frac{\phi_{j^*(\omega)}(f_\omega)}
       {\min_{k \in \Gl} \phik(f_\omega)}
  \;\geq\; 1 + \Omega(\rhovar).
\end{equation}
As $\rhovar \to 1$, the attribution ratio diverges: the first-mover
captures an increasingly dominant share of the group's total importance.
Hence no single model $f_\omega$ is $\gamma$-equitable for bounded
$\gamma$.

\medskip
\noindent\textbf{Part (ii): Equity enforcement destroys Stability.}

Suppose one attempts to restore equity by some mechanism that produces
$\phij \approx \phik$ for all $j, k \in \Gl$ (e.g., permutation-based
renormalization, or averaging within the group for a single model).
Consider what this implies for the ranking of features across $\Gl$.

For two independent seeds $\omega$ and $\omega'$, the first-movers
$j^*(\omega)$ and $j^*(\omega')$ are drawn approximately uniformly from
$\Gl$ (by DGP symmetry).  With probability $1 - 1/m$, they differ:
$j^*(\omega) \neq j^*(\omega')$.

\emph{Without equity enforcement}, the importance ranking within $\Gl$ is
determined by the first-mover identity.  When
$j^*(\omega) \neq j^*(\omega')$, the top-ranked feature within $\Gl$
differs between the two runs, and the Spearman correlation of the
full importance vectors satisfies
\begin{equation}
  \mathrm{Spearman}\bigl(\boldsymbol{\phi}(f_\omega),\;
    \boldsymbol{\phi}(f_{\omega'})\bigr)
  \;\leq\; 1 - \Omega\!\left(\frac{m}{P}\right),
\end{equation}
since at least $m - 1$ features in $\Gl$ receive reshuffled ranks.
As $\rhovar \to 1$, the concentration effect strengthens and the
rank disagreement within $\Gl$ becomes a dominant source of instability.

\emph{With equity enforcement}, the attribution gap within $\Gl$ is
compressed to near zero.  But the rank \emph{ordering} of near-tied
features is then determined by residual noise of order
$O(1/\sqrt{n})$ in the SHAP estimates, which is independent across
seeds.  The within-group ranking therefore becomes effectively random:
\begin{equation}
  \Exp\bigl[\mathrm{Spearman}_{\Gl}\bigl(
    \boldsymbol{\phi}(f_\omega),\;
    \boldsymbol{\phi}(f_{\omega'})
  \bigr)\bigr] \;\to\; 0
  \quad \text{as } \rhovar \to 1.
\end{equation}
In either case, stability degrades with $\rhovar$.

\medskip
\noindent\textbf{Combining both parts.}

A single sequential gradient-boosted model $f_\omega$ faces a dichotomy:
\begin{itemize}
  \item \emph{Accept the natural split pattern}: Equity is violated by
    $\Omega(\rhovar)$ (\Cref{lem:split-concentration,lem:shap-inherits}).
  \item \emph{Enforce equitable attributions}: The resulting near-ties
    make within-group rankings noise-dominated, and since the noise is
    independent across seeds, Stability is violated.
\end{itemize}
No finite $\totalT$, learning rate $\eta$, or tree depth $d$ can resolve
this, because the first-mover advantage in
\Cref{lem:split-concentration} holds for any $\totalT \geq 1$ and any
$\eta > 0$, and the symmetry-breaking in $h_1$ is an inherent
consequence of greedy sequential optimization.
\end{proof}


% ══════════════════════════════════════════════════════════════════════════════
\section{DASH Circumvention}
\label{sec:circumvention}
% ══════════════════════════════════════════════════════════════════════════════

\begin{corollary}[\dash{} circumvents the impossibility]
\label{cor:dash}
Let $f_{\omega_1}, \dots, f_{\omega_M}$ be $M$ independently trained
gradient-boosted models with i.i.d.\ seeds $\omega_1, \dots, \omega_M$.
Define the \dash{} consensus attribution as
\begin{equation}
  \bar{\phi}_j = \frac{1}{M} \sum_{i=1}^{M} \phij(f_{\omega_i}).
\end{equation}
Then:
\begin{enumerate}
  \item[(a)] \textbf{Equity in expectation.} By symmetry of the DGP and
    independence of the seeds, for any $j, k \in \Gl$,
    \begin{equation}
      \Exp[\bar{\phi}_j] = \Exp[\bar{\phi}_k].
    \end{equation}
    The first-mover identity $j^*(\omega_i)$ is approximately uniformly
    distributed over $\Gl$, so each feature serves as first-mover in
    approximately $M/m$ of the $M$ models.  The attribution concentration
    from \Cref{lem:shap-inherits} cancels in expectation.

  \item[(b)] \textbf{Stability via the law of large numbers.}
    By independence of $\{\omega_i\}_{i=1}^M$, the variance of the
    consensus attribution decays as
    \begin{equation}
      \Var(\bar{\phi}_j)
        = \frac{1}{M}\, \Var\bigl(\phij(f_{\omega_1})\bigr)
        = O(1/M).
    \end{equation}
    For sufficiently large $M$, two independent \dash{} runs (with
    different meta-seeds controlling the $M$ individual seeds) produce
    consensus vectors $\bar{\boldsymbol{\phi}}$ and
    $\bar{\boldsymbol{\phi}}'$ satisfying
    \begin{equation}
      \Exp\bigl[\mathrm{Spearman}\bigl(
        \bar{\boldsymbol{\phi}},\;
        \bar{\boldsymbol{\phi}}'
      \bigr)\bigr] \;\to\; 1
      \quad \text{as } M \to \infty.
    \end{equation}
\end{enumerate}
\dash{} thus achieves both Equity and Stability simultaneously by
replacing a single sequential model with an average over an independent
ensemble, breaking the conditions of \Cref{thm:impossibility}.
\end{corollary}

\begin{proof}
Part~(a) follows from linearity of expectation and the symmetry of the
DGP under permutation of features within each group $\Gl$.  Since the
seeds $\omega_i$ are i.i.d., the probability that $j^*(\omega_i) = j$
is $1/m$ for each $j \in \Gl$ and each $i$.  Hence
$\Exp[\phij(f_{\omega_i})] = \Exp[\phik(f_{\omega_i})]$ for all
$j, k \in \Gl$, and linearity gives
$\Exp[\bar{\phi}_j] = \Exp[\bar{\phi}_k]$.

Part~(b) follows from the standard law of large numbers.
The attributions $\phij(f_{\omega_1}), \dots, \phij(f_{\omega_M})$ are
i.i.d.\ with finite variance (TreeSHAP values are bounded for bounded
tree ensembles on bounded data).  Therefore
$\bar{\phi}_j \xrightarrow{\mathrm{a.s.}} \Exp[\phij(f_\omega)]$, and
the concentration inequality
$\Prob(\abs{\bar{\phi}_j - \Exp[\bar{\phi}_j]} > \epsilon)
  \leq \Var(\phij) / (M \epsilon^2)$
ensures that two independent \dash{} consensus vectors converge to the
same limit, yielding perfect rank agreement asymptotically.
\end{proof}


% ══════════════════════════════════════════════════════════════════════════════
\section{Discussion of Assumptions}
\label{sec:assumptions}
% ══════════════════════════════════════════════════════════════════════════════

\begin{remark}[Sequential structure is essential]
\label{rem:sequential}
The impossibility in \Cref{thm:impossibility} relies critically on the
\emph{sequential} nature of gradient boosting, where each tree $h_t$ is
fit to residuals $r_t = y - \sum_{s<t} h_s(x)$.  This creates a causal
chain: the first tree's split choices alter the residual landscape for all
subsequent trees, producing the cumulative first-mover advantage of
\Cref{lem:split-concentration}.

The theorem does \emph{not} apply to methods that train base learners
independently.  For example:
\begin{itemize}
  \item \textbf{Random forests} train each tree on a bootstrap sample with
    independent feature sub-sampling.  Since trees do not share residuals,
    there is no first-mover propagation: if tree $t$ splits on $j$, this
    does not affect tree $t'$'s feature selection.  The equity violation is
    therefore $O(1)$ rather than $O(\rhovar \cdot \totalT)$.
  \item \textbf{Bagged ensembles} similarly avoid the sequential dependency
    by construction.
\end{itemize}
However, gradient boosting's sequential structure is precisely what makes
it a powerful function approximator (by fitting residuals, it can represent
complex interactions that parallel ensembles miss).  The impossibility
result thus highlights a fundamental tension between the \emph{modeling
power} of sequential boosting and the \emph{interpretability desideratum}
of equitable feature attribution.
\end{remark}

\begin{remark}[Role of $\rhovar$]
\label{rem:rho}
The impossibility is most severe as $\rhovar \to 1$ but is present for
any $\rhovar > 0$.  At moderate correlations (e.g., $\rhovar = 0.5$), the
first-mover advantage exists but may be small enough that finite-sample
noise in SHAP estimates masks the equity violation.  The theorem
establishes that no principled guarantee can be made for a single model,
regardless of $\rhovar$; the practical severity scales with $\rhovar$.
\end{remark}

\begin{remark}[Finite-sample considerations]
\label{rem:finite}
The asymptotic statements ($\rhovar \to 1$, $M \to \infty$) are
idealizations.  In practice, \dash{} uses finite $M$ (typically 200
population models, 30 selected) and operates at finite $\rhovar$.
The corollary's convergence rate of $O(1/M)$ for variance provides
practical guidance: even $M = 30$ substantially reduces within-group
attribution variance compared to a single model.  The empirical results
in the main paper confirm that this finite-$M$ regime suffices to achieve
near-perfect equity on synthetic DGPs with $\rhovar \geq 0.9$.
\end{remark}

\end{document}
